% =================================================
% Рисунки к учебному комплекту
% Все координаты основной конфигурации согласованы.
% =================================================

% -------------------------------------------------
% 1. Исходная трапеция
% -------------------------------------------------
\newcommand{\trapezoidproblemfigure}{%
\begin{center}
\begin{tikzpicture}[scale=0.72,every node/.style={geo label}]
  \coordinate (B) at (0,0);
  \coordinate (C) at (8,0);
  \coordinate (N) at (3.31,9.615);
  \coordinate (A) at ($(B)!0.5!(N)$);
  \coordinate (D) at ($(C)!0.5!(N)$);
  \coordinate (M) at (4,4);

  \draw[geo main] (A)--(B)--(C)--(D)--cycle;
  \draw[geo radius] (A)--(M)--(D);

  \node[geo point,label={[geo label]above left:$A$}] at (A) {};
  \node[geo point,label={[geo label]below left:$B$}] at (B) {};
  \node[geo point,label={[geo label]below right:$C$}] at (C) {};
  \node[geo point,label={[geo label]above right:$D$}] at (D) {};
  \node[geo center point,label={[geo label]below:$M$}] at (M) {};

  \pic[geo right angle] {right angle=B--A--M};
  \pic[geo right angle] {right angle=M--D--C};
  \pic[geo angle,"$64^\circ$"{geo angle label},angle radius=8mm,angle eccentricity=1.32]
    {angle=D--C--B};
\end{tikzpicture}
\end{center}%
}

% -------------------------------------------------
% 2. Достроение трапеции до треугольника
% -------------------------------------------------
\newcommand{\hiddentrianglefigure}{%
\begin{center}
\begin{tikzpicture}[scale=0.55,every node/.style={geo label}]
  \coordinate (B) at (0,0);
  \coordinate (C) at (8,0);
  \coordinate (N) at (3.31,9.615);
  \coordinate (A) at ($(B)!0.5!(N)$);
  \coordinate (D) at ($(C)!0.5!(N)$);

  \draw[geo main] (B)--(N)--(C)--cycle;
  \draw[geo strong] (A)--(D);

  \node[geo point,label={[geo label]left:$A$}] at (A) {};
  \node[geo point,label={[geo label]below left:$B$}] at (B) {};
  \node[geo point,label={[geo label]below right:$C$}] at (C) {};
  \node[geo point,label={[geo label]right:$D$}] at (D) {};
  \node[geo point,label={[geo label]above:$N$}] at (N) {};
\end{tikzpicture}
\end{center}%
}

% -------------------------------------------------
% 3. Средняя линия: та же конфигурация и масштаб
% -------------------------------------------------
\newcommand{\midlinefigure}{%
\begin{center}
\begin{tikzpicture}[scale=0.66,every node/.style={geo label}]
  \coordinate (B) at (0,0);
  \coordinate (C) at (8,0);
  \coordinate (N) at (3.31,9.615);
  \coordinate (A) at ($(B)!0.5!(N)$);
  \coordinate (D) at ($(C)!0.5!(N)$);

  \draw[geo main] (B)--(N)--(C)--cycle;
  \draw[geo highlight] (A)--(D);

  \node[geo point,label={[geo label]left:$A$}] at (A) {};
  \node[geo point,label={[geo label]below left:$B$}] at (B) {};
  \node[geo point,label={[geo label]below right:$C$}] at (C) {};
  \node[geo point,label={[geo label]right:$D$}] at (D) {};
  \node[geo point,label={[geo label]above:$N$}] at (N) {};

  % Равные части стороны BN
  \draw[geo equal mark] ($(B)!0.58!(A)$) ++(-0.11,0.04) -- ++(0.22,-0.08);
  \draw[geo equal mark] ($(A)!0.42!(N)$) ++(-0.11,0.04) -- ++(0.22,-0.08);
  % Равные части стороны CN (двойные засечки)
  \draw[geo equal mark] ($(C)!0.58!(D)$) ++(-0.10,-0.05) -- ++(0.20,0.10);
  \draw[geo equal mark] ($(C)!0.58!(D)$) ++(-0.16,-0.08) -- ++(0.20,0.10);
  \draw[geo equal mark] ($(D)!0.42!(N)$) ++(-0.10,-0.05) -- ++(0.20,0.10);
  \draw[geo equal mark] ($(D)!0.42!(N)$) ++(-0.16,-0.08) -- ++(0.20,0.10);
\end{tikzpicture}
\end{center}%
}

% -------------------------------------------------
% 4. Серединный перпендикуляр
% -------------------------------------------------
\newcommand{\perpendicularbisectorfigure}{%
\begin{center}
\begin{tikzpicture}[scale=0.86,every node/.style={geo label}]
  \coordinate (B) at (0,0);
  \coordinate (N) at (6,0);
  \coordinate (A) at (3,0);
  \coordinate (M) at (3,3.1);

  \draw[geo main] (B)--(N);
  \draw[geo highlight] (A)--(M);
  \draw[geo radius] (M)--(B) (M)--(N);

  \node[geo point,label={[geo label]below:$A$}] at (A) {};
  \node[geo point,label={[geo label]below left:$B$}] at (B) {};
  \node[geo point,label={[geo label]below right:$N$}] at (N) {};
  \node[geo center point,label={[geo label]above:$M$}] at (M) {};

  \pic[geo right angle] {right angle=B--A--M};
  \draw[geo equal mark] (1.47,0.13)--(1.47,-0.13);
  \draw[geo equal mark] (4.53,0.13)--(4.53,-0.13);
\end{tikzpicture}
\end{center}%
}

% -------------------------------------------------
% 5. Центр описанной окружности
% -------------------------------------------------
\newcommand{\circumcenterfigure}{%
\begin{center}
\begin{tikzpicture}[scale=0.60,every node/.style={geo label}]
  \coordinate (B) at (0,0);
  \coordinate (C) at (8,0);
  \coordinate (N) at (3.31,9.615);
  \coordinate (A) at ($(B)!0.5!(N)$);
  \coordinate (D) at ($(C)!0.5!(N)$);
  \coordinate (M) at (4,4);

  \draw[geo circle] (M) circle[radius=5.657];
  \draw[geo main] (B)--(N)--(C)--cycle;
  \draw[geo auxiliary] (A)--(M) (D)--(M);
  \draw[geo radius] (M)--(B) (M)--(N) (M)--(C);

  \node[geo point,label={[geo label]left:$A$}] at (A) {};
  \node[geo point,label={[geo label]below left:$B$}] at (B) {};
  \node[geo point,label={[geo label]below right:$C$}] at (C) {};
  \node[geo point,label={[geo label]right:$D$}] at (D) {};
  \node[geo point,label={[geo label]above:$N$}] at (N) {};
  \node[geo center point,label={[geo label]below right:$M$}] at (M) {};

  \pic[geo right angle,angle radius=3.6mm] {right angle=B--A--M};
  \pic[geo right angle,angle radius=3.6mm] {right angle=M--D--C};
\end{tikzpicture}
\end{center}%
}

% -------------------------------------------------
% 6. Перпендикуляр ME к BC
% -------------------------------------------------
\newcommand{\perpendicularmefigure}{%
\begin{center}
\begin{tikzpicture}[scale=0.82,every node/.style={geo label}]
  \coordinate (B) at (0,0);
  \coordinate (C) at (8,0);
  \coordinate (M) at (4,3.5);
  \coordinate (E) at (4,0);

  \draw[geo main] (B)--(C);
  \draw[geo radius] (M)--(B) (M)--(C);
  \draw[geo highlight] (M)--(E);

  \node[geo point,label={[geo label]below left:$B$}] at (B) {};
  \node[geo point,label={[geo label]below right:$C$}] at (C) {};
  \node[geo center point,label={[geo label]above:$M$}] at (M) {};
  \node[geo point,label={[geo label]below:$E$}] at (E) {};

  \pic[geo right angle] {right angle=B--E--M};
  \draw[geo equal mark] (2,0.13)--(2,-0.13);
  \draw[geo equal mark] (6,0.13)--(6,-0.13);
\end{tikzpicture}
\end{center}%
}

% -------------------------------------------------
% 7. Прямоугольный равнобедренный треугольник BME
% -------------------------------------------------
\newcommand{\trianglebmefigure}{%
\begin{center}
\begin{tikzpicture}[scale=0.82,every node/.style={geo label}]
  \coordinate (B) at (0,0);
  \coordinate (E) at (4,0);
  \coordinate (M) at (4,4);

  \draw[geo main] (B)--(E)--(M)--cycle;
  \node[geo point,label={[geo label]below left:$B$}] at (B) {};
  \node[geo point,label={[geo label]below right:$E$}] at (E) {};
  \node[geo center point,label={[geo label]above:$M$}] at (M) {};

  \pic[geo right angle] {right angle=B--E--M};
  \pic[geo angle,"$45^\circ$"{geo angle label},angle radius=8mm,angle eccentricity=1.28]
    {angle=E--B--M};
  \node[below] at ($(B)!0.5!(E)$) {$a$};
  \node[right] at ($(E)!0.5!(M)$) {$a$};
\end{tikzpicture}
\end{center}%
}

% -------------------------------------------------
% 8. Центральный и вписанный углы
% -------------------------------------------------
\newcommand{\centralinscribedfigure}{%
\begin{center}
\begin{tikzpicture}[scale=0.60,every node/.style={geo label}]
  \coordinate (B) at (0,0);
  \coordinate (C) at (8,0);
  \coordinate (N) at (3.31,9.615);
  \coordinate (M) at (4,4);

  \draw[geo circle] (M) circle[radius=5.657];
  \draw[geo main] (B)--(C)--(N)--cycle;
  \draw[geo radius] (M)--(B) (M)--(N);

  \node[geo point,label={[geo label]below left:$B$}] at (B) {};
  \node[geo point,label={[geo label]below right:$C$}] at (C) {};
  \node[geo point,label={[geo label]above:$N$}] at (N) {};
  \node[geo center point,label={[geo label]below right:$M$}] at (M) {};

  \pic[geo angle,"$64^\circ$"{geo angle label},angle radius=8mm,angle eccentricity=1.35]
    {angle=B--C--N};
  \pic[geo angle,"$128^\circ$"{geo angle label},angle radius=10mm,angle eccentricity=1.24]
    {angle=B--M--N};
\end{tikzpicture}
\end{center}%
}

% -------------------------------------------------
% 9. Равнобедренный треугольник BMN
% -------------------------------------------------
\newcommand{\trianglebmnfigure}{%
\begin{center}
\begin{tikzpicture}[scale=0.68,every node/.style={geo label}]
  \coordinate (B) at (0,0);
  \coordinate (M) at (4,4);
  \coordinate (N) at (3.31,9.615);

  \draw[geo main] (B)--(M)--(N)--cycle;
  \node[geo point,label={[geo label]below left:$B$}] at (B) {};
  \node[geo center point,label={[geo label]right:$M$}] at (M) {};
  \node[geo point,label={[geo label]above:$N$}] at (N) {};

  \pic[geo angle,"$128^\circ$"{geo angle label},angle radius=7.5mm,angle eccentricity=1.30]
    {angle=B--M--N};
  \pic[geo angle,"$26^\circ$"{geo angle label},angle radius=7.5mm,angle eccentricity=1.28]
    {angle=M--B--N};
  \pic[geo angle,"$26^\circ$"{geo angle label},angle radius=7mm,angle eccentricity=1.28]
    {angle=B--N--M};

  \draw[geo equal mark] ($(M)!0.48!(B)$) ++(-0.10,0.10)--++(0.20,-0.20);
  \draw[geo equal mark] ($(M)!0.48!(N)$) ++(-0.10,-0.02)--++(0.20,0.04);
\end{tikzpicture}
\end{center}%
}

% -------------------------------------------------
% 10. Итоговое сложение углов
% -------------------------------------------------
\newcommand{\anglesumfigure}{%
\begin{center}
\begin{tikzpicture}[scale=0.70,every node/.style={geo label}]
  \coordinate (B) at (0,0);
  \coordinate (C) at (8,0);
  \coordinate (N) at (3.31,9.615);
  \coordinate (M) at (4,4);

  \draw[geo main] (B)--(N);
  \draw[geo main] (B)--(C);
  \draw[geo highlight] (B)--(M);

  \node[geo point,label={[geo label]below left:$B$}] at (B) {};
  \node[geo point,label={[geo label]below right:$C$}] at (C) {};
  \node[geo center point,label={[geo label]right:$M$}] at (M) {};
  \node[geo point,label={[geo label]above:$N$}] at (N) {};

  \pic[geo angle,"$26^\circ$"{geo angle label},angle radius=9mm,angle eccentricity=1.28]
    {angle=N--B--M};
  \pic[geo angle,"$45^\circ$"{geo angle label},angle radius=15mm,angle eccentricity=1.18]
    {angle=M--B--C};
  \pic[draw=Primary,line width=0.75pt,"$71^\circ$"{geo angle label},angle radius=22mm,angle eccentricity=1.10]
    {angle=N--B--C};
\end{tikzpicture}
\end{center}%
}

% -------------------------------------------------
% Дополнительная схема для ошибки N,M,E
% -------------------------------------------------
\newcommand{\collinearityerrorfigure}{%
\begin{center}
\begin{tikzpicture}[scale=0.66,every node/.style={geo label}]
  \coordinate (B) at (0,0);
  \coordinate (C) at (8,0);
  \coordinate (N) at (3.31,9.615);
  \coordinate (M) at (4,4);
  \coordinate (E) at (4,0);
  \draw[geo main] (B)--(N)--(C)--cycle;
  \draw[geo highlight] (M)--(E);
  \draw[geo construction] (N)--(M);
  \node[geo point,label={[geo label]below left:$B$}] at (B) {};
  \node[geo point,label={[geo label]below right:$C$}] at (C) {};
  \node[geo point,label={[geo label]above:$N$}] at (N) {};
  \node[geo center point,label={[geo label]right:$M$}] at (M) {};
  \node[geo point,label={[geo label]below:$E$}] at (E) {};
  \pic[geo right angle] {right angle=B--E--M};
  \node[font=\bfseries\sffamily\Large,text=Accent] at (5.85,5.7) {$\times$};
\end{tikzpicture}
\end{center}%
}
